\clearpage

\chapter{Methodology}

In this chapter, we outline the methodology used to develop and evaluate the procedural
quest generation system for RPGs using LLMs. The goal of this research is to create an
automated framework that can generate coherent, engaging, and contextually appropriate
quests for RPGs, leveraging the power of LLMs and hybrid architectures. This chapter
describes the step-by-step procedure, from data collection and pre-processing to model
training, fine-tuning, and evaluation. The methodology is designed to ensure that the
generated quests are of high quality, both from a computational and a user experience
perspective.

We begin by discussing the dataset collection and preparation, including the specifics
of the quest data used to train the models. Next, we detail the architecture of the LLM-based
models, the fine-tuning procedures, and the use of LoRA adapters to optimize model
performance for quest generation. Following that, we delve into the evaluation metrics
employed to assess the generated quests, explaining both quantitative and qualitative
methods such as automatic evaluation metrics and human evaluation. Finally, we outline
the experimental setup and the procedures followed during model training, validation,
and testing, with a focus on how the methodology ensures the generation of diverse and
contextually rich quests for RPGs.

\section{Workflow for Procedural Quest Generation}

This section provides an end-to-end overview of the PQG workflow implemented in this
study. The workflow integrates multiple components, from dataset preparation to model
fine-tuning and evaluation, to produce coherent, context-aware, and narratively rich quests
for RPGs. Each phase of the workflow has been designed to ensure the robustness,
efficiency, and adaptability of the quest generation system. The workflow is structured to
align with current best practices in NLG using instruction-tuned LLMs, while addressing
the unique challenges of interactive narrative construction in games.

The system emphasizes modularity across all stages, allowing rapid experimentation
with models, preprocessing, and evaluation strategies. From curated, annotated quest
datasets to prompt-based generation with lightweight fine-tuning, each component supports
a cohesive and flexible pipeline. Quantization and parameter-efficient tuning (e.g.,
LoRA) enable deployment in constrained environments without compromising quality.
Generation and evaluation setups are optimized to capture both lexical accuracy and
semantic depth. The complete workflow is represented by the flow diagram depicted in Figure~\ref{fig:pqg-workflow}.

\begin{figure}[t]
  \centering
  \begin{tikzpicture}[
    node distance=1cm,
    every node/.style={font=\scriptsize},
    step/.style={rectangle, draw=black, rounded corners, minimum width=4.5cm, minimum height=1.2cm, text centered, fill=#1!30, drop shadow},
    desc/.style={align=left, anchor=west, font=\footnotesize, text width=8cm},
    arrow/.style={-{Latex[length=2mm]}, thick}
  ]
    \node[step=blue] (data) {Dataset Collection};
    \node[step=cyan, below=of data] (format) {Tagging \& Formatting};
    \node[step=teal, below=of format] (prompt) {Prompt Construction};
    \node[step=lime, below=of prompt] (model) {Model Selection};
    \node[step=green, below=of model] (quant) {Quantization Setup};
    \node[step=olive, below=of quant] (lora) {LoRA Adapter Integration};
    \node[step=orange, below=of lora] (train) {Training Pipeline};
    \node[step=yellow!80!brown, below=of train] (gen) {Generation Setup};
    \node[step=red!60!brown, below=of gen] (eval) {Evaluation Metrics};

    \node[desc, right=1.5cm of data] {Collecting domain-specific corpora consisting of fantasy and role-playing game quests to serve as training data for the generation pipeline.};
    \node[desc, right=1.5cm of format] {Tagging entities and narrative components using structured markers (e.g., \texttt{<|begin\_objective|>}, \texttt{<|begin\_tasks|>}) and formatting them into model-compatible input sequences.};
    \node[desc, right=1.5cm of prompt] {Creating template-based prompts that incorporate domain-specific tokens to condition the model towards generating structured quest narratives.};
    \node[desc, right=1.5cm of model] {Selecting pre-trained language models such as GPT-2 or LLaMA that support integration with LoRA adapters and quantized fine-tuning techniques.};
    \node[desc, right=1.5cm of quant] {Utilizing 4-bit quantization via QLoRA to significantly reduce memory usage, enabling training on consumer-grade hardware while preserving model performance.};
    \node[desc, right=1.5cm of lora] {Embedding LoRA adapters within attention blocks to perform parameter-efficient fine-tuning with reduced computational overhead.};
    \node[desc, right=1.5cm of train] {Executing the training process with integrated checkpointing and diagnostic tools to monitor convergence via metrics such as training loss and gradient norms.};
    \node[desc, right=1.5cm of gen] {Setting generation hyperparameters like temperature, top-k, and top-p to balance creativity and coherence in the produced quest outputs.};
    \node[desc, right=1.5cm of eval] {Evaluating generated quests using automated metrics (BLEU, METEOR, BERTScore) alongside manual inspection to assess lexical and semantic quality.};

    \foreach \i/\j in {data/format, format/prompt, prompt/model, model/quant,quant/lora, lora/train, train/gen, gen/eval}
      \draw[arrow] (\i) -- (\j);
  \end{tikzpicture}
  \caption{Flow diagram of the procedural quest generation system. It illustrates the modular pipeline used for model training and evaluation.}
  \label{fig:pqg-workflow}
\end{figure}
\FloatBarrier

\subsection{Dataset Collection and Formatting}

The process begins with dataset preparation, a foundational step critical for training
effective quest generation models. Due to the scarcity of structured, publicly available
quest datasets, data is collected from a variety of sources, including open-source RPG
scripts (e.g., \textit{Skyrim}~\cite{theelderscrollsvskyrim} and \textit{Baldur's Gate}~\cite{baldursgate}), community-maintained wikis, and synthetic
examples generated using quest templates.

The primary dataset used in this study is the V{\"a}rtinen et al.~\cite{vartinen2022generating} dataset, which
consists of 978 manually annotated quests extracted from a diverse set of RPG titles,
including \textit{Minecraft}~\cite{minecraft} and \textit{Torchlight II}~\cite{torchlightii}. Each quest in this dataset is represented
in a structured XML-inspired format, comprising tags for key quest elements including
\texttt{<|begin\_objective|>}, \texttt{<|begin\_tasks|>}, and \texttt{<|begin\_quest\_giver|>}. This format
maintains the meaningful connections between quest elements while enabling efficient
structured data processing for subsequent tasks.

To enhance the dataset's scope and represent a broader variety of game narratives, this
study appends 156 additional quests from the \textit{Fallout series}~\cite{fallout1,fallout2}. These quests are extracted
manually from community-maintained sources and restructured to align with the
XML-like format used by Vartinen et al.~\cite{vartinen2022generating}, ensuring consistency and interoperability.

After consolidation, the dataset underwent preprocessing to resolve inconsistencies
and segment each quest into well-defined narrative components. The structured quests
are then transformed into an instruction-style prompt-response format to suit instruction-tuned
language models. A sample of the formatted quests and the prompt-response
formatting can be seen in Appendix~\ref{appendix:quest-dataset}.

\subsection{Model Selection and Evaluation}

Following dataset construction, the workflow progressed to model selection, quantization,
and fine-tuning. Multiple pretrained language models are evaluated, including
GPT-2 variants, TinyLLaMA, and LLaMA 3.2, to compare their performance in PQG.
These models are chosen for their varying parameter scales, open availability, and relevance
to instruction-based generation tasks.

To enable training on resource-constrained, consumer-grade hardware, all models are
quantized to 4-bit precision using QLoRA (Quantized Low-Rank Adaptation)~\cite{dettmers2023qlora}. This
method significantly reduces memory usage and computational overhead while maintaining
most of the model's generative capabilities. LoRA is used as the primary PEFT
method, enabling efficient updates through low-rank projections without modifying the
entire model architecture~\cite{peft,hu2022lora}.

Fine-tuning is conducted using the XML-inspired, instruction-annotated dataset, converted
into prompt-response pairs. Although true mixed-precision training is not supported
due to backend limitations, float16 precision is employed where possible to reduce
memory footprint. Gradient checkpointing is utilized to further optimize memory
usage during backpropagation, and early stopping is incorporated to prevent overfitting based
on validation loss trends. The training objective centered on minimizing cross-entropy
loss between model outputs and target quest responses, enhancing the model's ability to
produce coherent, context-aware, and narratively rich quests.

Once fine-tuning is complete, model evaluation is conducted using a combination of
automated metrics and comparative performance analysis. Automated evaluation involved
standard NLG metrics, including BLEU, METEOR, and BERTScore, to assess
lexical overlap, semantic alignment, and fluency between generated and reference quest
narratives. In addition to these, perplexity is used as an intrinsic measure of the model's
generative confidence, with lower perplexity indicating better language modeling performance.

\subsection{Framework Scope and Limitations}

While deployment and runtime integration are not the primary objectives of this study,
the system has been developed as a procedural framework for advancing quest generation
in low-resource environments. The emphasis is placed on designing an adaptable pipeline
that demonstrates the feasibility of generating context-aware and narratively coherent
quests using quantized language models and PEFT methods. Rather than targeting
production-ready deployment, the goal is to provide a robust experimental foundation
that can be extended with larger models, richer datasets, and more advanced training
configurations.

\section{Dataset Description and Preprocessing}

This section introduces the dataset compiled for training and evaluating LLMs in the task
of PCG. The construction of the dataset involves aggregating quest data from diverse
RPG sources, standardizing the structure into a uniform schema, and formatting it for
compatibility with generative language models.

Given the narrative and structural complexity of quests in RPGs, it is essential to
represent each quest in a machine-readable yet semantically rich format that preserves
critical gameplay elements such as objectives, subtasks, quest givers, rewards, and supporting
entities. This requires a preprocessing pipeline capable of handling heterogeneous
data from multiple games, some of which originate from structured datasets while others
are manually curated and normalized with the aid of LLMs.

The final dataset, referred to as the QUEST Dataset, captures a wide range of quest
archetypes and contextual variations across both high-fantasy and post-apocalyptic game
settings. It combines data from multiple titles, including games from the \textit{Elder Scrolls,
Fallout, and Baldur's Gate franchises}~\cite{theelderscrollsivoblivion,theelderscrollsvskyrim,fallout1,fallout2,baldursgate,baldursgate2shadowsofamn}, and ensures each entry adheres to a unified
tagging convention that facilitates token-level alignment during training.

To prepare the dataset for sequence modeling tasks, each quest instance is transformed
into a text sequence using an XML-inspired prompt structure. This formatting approach
allows for consistent parsing, enhances controllability during generation, and supports
both conditional and unconditional sampling in downstream experiments. The details of
data collection, game-wise contributions, and the exact schema used for formatting are
described in the following subsections.

\subsection{Source and Specification of Dataset}

\begin{table}[t]
  \centering
  \scriptsize
  \renewcommand{\arraystretch}{1.3}
  \begin{tabularx}{0.95\textwidth}{
    >{\raggedright\arraybackslash}p{6cm}
    >{\raggedright\arraybackslash}X
    >{\centering\arraybackslash}X
    >{\centering\arraybackslash}X
  }
    \toprule
    \textbf{Game Title} & \textbf{Source} & \textbf{Number of Quests} & \textbf{Percentage} \\
    \midrule
    Baldur's Gate~\cite{baldursgate} & V{\"a}rtinen et al.~\cite{vartinen2022generating} & 100 & 8.8183\% \\
    Baldur's Gate II: Shadows of Amn~\cite{baldursgate2shadowsofamn} & V{\"a}rtinen et al.~\cite{vartinen2022generating} & 94 & 8.2892\% \\
    The Elder Scrolls IV: Oblivion~\cite{theelderscrollsivoblivion} & V{\"a}rtinen et al.~\cite{vartinen2022generating} & 215 & 18.9594\% \\
    The Elder Scrolls V: Skyrim~\cite{theelderscrollsvskyrim} & V{\"a}rtinen et al.~\cite{vartinen2022generating} & 389 & 34.3034\% \\
    Minecraft~\cite{minecraft} & V{\"a}rtinen et al.~\cite{vartinen2022generating} & 100 & 8.8183\% \\
    Torchlight II~\cite{torchlightii} & V{\"a}rtinen et al.~\cite{vartinen2022generating} & 80 & 7.0547\% \\
    Fallout~\cite{fallout1} & Nukapedia~\cite{fallout1quests} & 58 & 5.1146\% \\
    Fallout 2~\cite{fallout2} & Nukapedia~\cite{fallout2quests} & 98 & 8.6420\% \\
    \midrule
    Total & \empty & 1134 & 99.9999\% \\
    \bottomrule
  \end{tabularx}
  \caption{Distribution of collected quest data categorized by the game title}
\end{table}

\begin{figure}[t]
  \begin{subfigure}[H]{0.3\textwidth}
    \centering
    \includegraphics[width=\textwidth,height=6.5cm]{fallout.jpg}
  \end{subfigure}
  \hfil
  \begin{subfigure}[H]{0.3\textwidth}
    \centering
    \includegraphics[width=\textwidth,height=6.5cm]{baldurs_gate_ii.png}
  \end{subfigure}
  \hfil
  \centering
  \begin{subfigure}[H]{0.3\textwidth}
    \centering
    \includegraphics[width=\textwidth,height=6.5cm]{skyrim.png}
  \end{subfigure}
  \caption{Cover arts for selected RPGs used in the QUEST dataset~\cite{fallout1,baldursgate2shadowsofamn,theelderscrollsvskyrim}}
\end{figure}

The QUEST dataset is constructed by aggregating quest data from a curated collection
of eight popular RPGs. These include titles from both Western and sandbox RPG traditions,
ensuring diversity in quest structure, theme, and complexity. The primary source of
this dataset is the open-access quest corpus published by V{\"a}rtinen et al.~\cite{vartinen2022generating}, which
includes quests from six mainstream RPGs. To enhance the dataset's size and narrative
richness, additional quests from \textit{Fallout}~\cite{fallout1} and \textit{Fallout 2}~\cite{fallout2} were manually collected
and structured using a semi-automated method involving an LLM assistant.

Each quest entry is composed of a set of core and optional fields that reflect the semantic
elements typically found in narrative design: goals, tasks, characters, environments,
and consequences. The data schema was extended slightly from the original dataset to
accommodate additional information present in the manually added entries (e.g., character
motivations, groups, and auxiliary items). These data fields are clearly defined in
Appendix~\ref{appendix:quest-dataset}.

The inclusion of quests from \textit{Fallout}~\cite{fallout1} and \textit{Fallout 2}~\cite{fallout2} expanded the original dataset
by 156 unique entries, enhancing both the size and variety of the corpus. Among all the titles
represented, \textit{The Elder Scrolls V: Skyrim}~\cite{theelderscrollsvskyrim} contributes the
largest share of quests, accounting for approximately 34.30\% of the total, whereas \textit{Fallout}~\cite{fallout1} contributes the fewest, at 5.11\%.

To support richer expression during generation, the dataset retains several sparsely
populated or “rare” columns—fields that are not consistently filled across all game sources
but are preserved for flexibility in representing optional or game-specific quest attributes.
As a result, the compiled dataset offers a comprehensive and heterogeneous benchmark
for evaluating the narrative coherence, structural consistency, and alignment with game-world
logic in procedurally generated quests.

\subsection{Tagging and Formatting}

To enable the dataset's compatibility with autoregressive LLMs, each quest entry is serialized
into a structured, tag-based format resembling XML. This formatting step ensures
consistency, simplifies prompt construction, and aligns with the model's tokenization
strategies during both training and inference.

The formatting approach employs custom tags that explicitly delineate each quest
field. Mandatory components—such as \texttt{<|begin\_objective|>}, \texttt{<|begin\_tasks|>}, and \texttt{<|begin\_quest\_giver|>}—are always included to guarantee minimum structure, while
optional fields are conditionally inserted based on content availability. This design accommodates
the dataset's heterogeneous nature and supports the generation of partial or
context-dependent quests.

\begin{lstlisting}[
  language=xml,
  caption={Snippet of the XML-like quest formatting template},
  xleftmargin=5.5cm,
]
<|begin_quest|>
<|begin_objective|>
[Objective goes here]
<|end_objective|>
<|begin_tasks|>
[Task 1]
[Task 2]
<|end_tasks|>
<|begin_quest_giver|>
[Name]: [Description]
<|end_quest_giver|>
...
<|end_quest|>
\end{lstlisting}

To support this transformation, a custom Python function (\texttt{format\_quest\_as\_xml})
was developed. It parses each quest dictionary, extracts structured elements (like \texttt{rewards},
\texttt{locations}, or \texttt{enemies}), and generates a well-formatted string using a consistent template.
Nested and list-based fields are joined with line breaks for model readability and
separation of elements. The complete tagging and formatting script is available in the
project repository\footnote{\url{https://github.com/DracoY-code/quest-generation}}.

This structured formatting approach plays a critical role in preparing the dataset for
model training. It enforces uniform tag usage, ensuring that all quest entries—regardless
of their original source—follow a standardized input format. The design also offers flexibility,
allowing non-essential fields to be selectively included or omitted without affecting
structural integrity.

Additionally, it facilitates controlled prompt construction, enabling fine-tuned exposure
of specific quest attributes during training and evaluation. By converting diverse,
game-specific data into a consistent and model-compatible format, this step serves as a
vital bridge between raw input and effective language model utilization. Further specifications
for the formatting procedure can be seen in Appendix~\ref{appendix:quest-dataset}.

\section{Model Configuration and Training}

This section describes the architecture, configuration, and training procedures used to
fine-tune language models for the task of PQG. Given the size and complexity of modern
LLMs, this phase incorporates efficient training strategies such as quantization and LoRA~\cite{hu2022lora} to
make the process more computationally tractable and memory efficient.

\subsection{Quantization Workflow}

\begin{table}[t]
  \centering
  \scriptsize
  \renewcommand{\arraystretch}{1.3}
  \begin{tabularx}{0.95\textwidth}{
    >{\raggedright\arraybackslash}p{5cm}
    >{\centering\arraybackslash}X
  }
    \toprule
    \textbf{Parameter} & \textbf{Value} \\
    \midrule
    Quantization Type & 4-bit NormalFloat (NF4) \\
    Compute Data Type & \texttt{bfloat16} / \texttt{float16} \\
    Double Quantization & Enabled \\
    Library Used & \texttt{bitsandbytes} with Hugging Face Transformers \\
    \bottomrule
  \end{tabularx}
  \caption{Quantization configurations used in model compression for training}
  \label{table:quant-config}
\end{table}

To reduce memory usage and accelerate training, model weights were quantized using
the BitsAndBytes\footnote{\url{https://github.com/bitsandbytes-foundation/bitsandbytes}} library. This allowed loading larger models into memory-constrained
environments while preserving inference quality. The quantization configuration is based
on the 4-bit NormalFloat (NF4) format with double quantization and computation performed
in 16-bit precision (bfloat16 or float16).

\begin{lstlisting}[
  language=python,
  caption={Quantization configuration using BitsAndBytes with 4-bit NF4 precision},
  xleftmargin=3cm,
]
bnb_config = BitsAndBytesConfig(
    load_in_4bit=True,
    bnb_4bit_quant_type="nf4",
    bnb_4bit_use_double_quant=True,
    bnb_4bit_compute_dtype=torch.float16,
)
\end{lstlisting}

This setup, as seen in Table~\ref{table:quant-config}, significantly reduces the GPU VRAM footprint
without degrading downstream performance, enabling multi-billion parameter models to
be trained on consumer-grade hardware or limited cloud instances.

\subsection{LoRA Adapter Integration}

\begin{table}[t]
  \centering
  \scriptsize
  \renewcommand{\arraystretch}{1.3}
  \begin{tabularx}{0.95\textwidth}{
    >{\raggedright\arraybackslash}p{5cm}
    >{\centering\arraybackslash}X
  }
    \toprule
    \textbf{Parameter} & \textbf{Value} \\
    \midrule
    Target Modules & Attention projection layers (\texttt{q\_proj}, \texttt{v\_proj}, \texttt{o\_proj}) \\
    Rank ($r$) & 4 \\
    Alpha ($\alpha$) & 8 \\
    Dropout & 0.0 \\
    Bias & None \\
    Library Used & \texttt{PEFT} (Parameter-Efficient Fine-Tuning) \\
    \bottomrule
  \end{tabularx}
  \caption{Configuration of LoRA adapter modules integrated during fine-tuning}
  \label{table:lora-adapter}
\end{table}

To further reduce the number of trainable parameters, Low-Rank Adaptation (LoRA)~\cite{hu2022lora}
modules were employed. LoRA injects trainable low-rank matrices into the attention
layers of transformer models, allowing efficient adaptation without modifying the base
model weights.

LoRA was selected for its efficiency and flexibility in fine-tuning LLMs. It enables
parameter-efficient training by updating only a small subset of parameters—typically less
than 1\% of the model's total—thereby significantly reducing both computational and
storage overhead. Its modular nature allows adapters to be added or removed without
altering the base model weights, facilitating plug-and-play experimentation across
different configurations.

\begin{lstlisting}[
  language=python,
  caption={LoRA configuration for causal language modeling using PEFT},
  xleftmargin=2.5cm,
]
lora_config = LoraConfig(
    r=4,
    lora_alpha=8,
    lora_dropout=0.0,
    target_modules=["q_proj", "v_proj", "o_proj"],
    bias="none",
    task_type="CAUSAL_LM",
)
\end{lstlisting}

Moreover, LoRA adapters support better generalization by mitigating catastrophic forgetting,
preserving the base model's pre-trained knowledge while enabling effective learning
of task-specific objectives. This configuration is implemented using Hugging Face's
\texttt{peft}\footnote{\url{https://github.com/huggingface/peft}} library, which provides seamless integration of LoRA within transformer-based
models. The configuration details of the LoRA adapter modules are listed in Table~\ref{table:lora-adapter}.

\subsection{Training Pipeline Setup}

\begin{table}[t]
  \centering
  \scriptsize
  \renewcommand{\arraystretch}{1.3}
  \begin{tabularx}{0.95\textwidth}{
    >{\raggedright\arraybackslash}p{5cm}
    >{\centering\arraybackslash}X
  }
    \toprule
    \textbf{Feature} & \textbf{Details} \\
    \midrule
    Frameworks & \texttt{transformers}, \texttt{datasets}, \texttt{accelerate}, \texttt{peft} \\
    Trainer & \texttt{transformers.Trainer} API \\
    Tokenization & Custom tokenizer with special XML-like tokens (e.g., \texttt{<|begin\_objective|>}) \\
    Loss function & Cross-entropy \\
    Evaluation & Perplexity, BLEU, METEOR, BERTScore on a held-out test split \\
    Logging & Optional integration with TensorBoard \\
    \bottomrule
  \end{tabularx}
  \caption{Features and components included in the training pipeline}
  \label{table:train-pipeline}
\end{table}

This training pipeline uses Hugging Face's \texttt{transformers}\footnote{\url{https://github.com/huggingface/transformers}}, \texttt{accelerate}\footnote{\url{https://github.com/huggingface/accelerate}}, and \texttt{datasets}\footnote{\url{https://github.com/huggingface/datasets}} frameworks.
The tokenized dataset is streamed into the model in batched form with support for gradient
accumulation, mixed precision, and distributed training.

This training loop supports both standalone LLM fine-tuning and adapter-based modular
fine-tuning and can be reused across multiple datasets or model architectures with
minimal reconfiguration. The training configurations are listed in Table~\ref{table:train-pipeline}.

\subsection{Text Generation Settings}
\label{section:text-generation-settings}

\begin{table}[t]
  \centering
  \scriptsize
  \renewcommand{\arraystretch}{1.3}
  \begin{tabularx}{0.95\textwidth}{
    >{\raggedright\arraybackslash}p{5cm}
    >{\centering\arraybackslash}X
  }
    \toprule
    \textbf{Parameter} & \textbf{Description} \\
    \midrule
    Prompt structure & XML-like tags (e.g., \texttt{<|begin\_objective|>}) \\
    Max sequence length & 512 tokens \\
    Temperature & 0.7 \\
    Top-k sampling & 50 \\
    Top-p sampling & 0.9 \\
    Repetition penalty & 1.2 \\
    \bottomrule
  \end{tabularx}
  \caption{Settings used for text generation tasks}
  \label{tab:generation-settings}
\end{table}

This subsection details the configuration and hyperparameters applied during the text
generation phase, which critically affect the quality, diversity, and coherence of the generated
quests. The generation parameters were carefully chosen to balance creativity with
structural fidelity, considering the specific characteristics of each model architecture, such
as GPT-2 and its fine-tuned variants.

These parameters listed in Table~\ref{tab:generation-settings} were designed to generate high-quality, diverse quest narratives whiles
maintaining structural consistency, ensuring that the generated outputs could be effectively
evaluated using both automated metrics and human assessments.

\section{Evaluation Strategy}

To assess the quality of generated quests, a combination of quantitative and qualitative
evaluation methods is employed. These include standard language modeling metrics like
perplexity, reference-based n-gram overlap metrics such as BLEU and METEOR, semantic
similarity scoring through BERTScore, and a qualitative analysis of sample outputs to
assess structural and narrative coherence. Each metric offers complementary insight into
different aspects of the generation quality: fluency, relevance, semantic fidelity, and task-specific
performance.

\subsection{Perplexity}

Perplexity\footnote{\url{https://huggingface.co/spaces/evaluate-metric/perplexity}} is a widely used intrinsic metric for evaluating language models. It quantifies
how well a probability model predicts a given sequence. Formally, the perplexity of a
language model \( M \) on a sequence of tokens \( w_1, w_2, ..., w_N \) is defined as:

\begin{equation}
  \text{Perplexity}(M) = \exp\left(-\frac{1}{N} \sum_{i=1}^{N} \log P(w_i \mid w_{<i})\right)
\end{equation}

Lower perplexity indicates that the model assigns higher probability to the actual
next tokens in the sequence, suggesting better fluency and alignment with the target
distribution.

\subsection{BLEU}

The Bilingual Evaluation Understudy\footnote{\url{https://huggingface.co/spaces/evaluate-metric/bleu}} (BLEU) score is a reference-based metric originally
developed for machine translation. It measures the n-gram overlap between the generated
output and one or more reference texts. The BLEU score is calculated using precision
over n-grams and includes a brevity penalty to penalize overly short outputs. For up to
4-grams, the score is given by:

\begin{equation}
  \text{BLEU} = \text{BP} \cdot \exp\left( \sum_{n=1}^{4} w_n \log p_n \right)
\end{equation}

where \( p_n \) is the modified precision for n-grams, \( w_n \) is the weight (typically \( \frac{1}{4} \)), and BP
is the brevity penalty, defined as:

\begin{equation}
  \text{BP} = 
    \begin{cases}
    1 & \text{if } c > r \\
    \exp\left(1 - \frac{r}{c}\right) & \text{if } c \leq r
    \end{cases}
\end{equation}

where \( c \) is the length of the candidate and \( r \) is the length of the reference.

\subsection{METEOR}

METEOR\footnote{\url{https://huggingface.co/spaces/evaluate-metric/meteor}} (Metric for Evaluation of Translation with Explicit ORdering) is another
reference-based metric that improves upon BLEU by incorporating stemming, synonymy
matching, and flexible alignment. It computes a harmonic mean of unigram precision and
recall, with a higher weight on recall. It also applies a fragmentation penalty to discourage
disordered alignments. The overall score is:

\begin{equation}
  \text{METEOR} = F_{\text{mean}} \cdot (1 - Penalty)
\end{equation}

and

\begin{equation}
  F_{\text{mean}} = \frac{10 \cdot P \cdot R}{R + 9P}
\end{equation}

where \( P \) and \( R \) are unigram precision and recall, respectively.

\subsection{BERTScore}

BERTScore\footnote{\url{https://huggingface.co/spaces/evaluate-metric/bertscore}} evaluates semantic similarity by leveraging contextual embeddings from
pre-trained BERT-like models. Instead of relying on exact n-gram overlap, it compares
each token in the candidate with the most similar token in the reference using cosine
similarity. The final score is a combination of precision, recall, and F1 computed over
these similarity scores:

\begin{equation}
  \text{BERTScore}_{F1} = \frac{2 \cdot P \cdot R}{P + R}
\end{equation}

where precision \( P \) and recall \( R \) are calculated using cosine similarity between contextualized
embeddings of matched tokens.

\subsection{Qualitative Analysis}

\begin{table}[t]
  \centering
  \scriptsize
  \renewcommand{\arraystretch}{1.3}
  \begin{tabularx}{0.95\textwidth}{
    >{\raggedright\arraybackslash}p{5cm}
    >{\raggedright\arraybackslash}X
  }
    \toprule
    \textbf{Criterion} & \textbf{Description} \\
    \midrule
    Narrative Coherence & Logical progression and internal consistency of the quest storyline \\
    Goal Clarity & Clear articulation of objectives and task sequences \\
    World Consistency & Adherence to the lore and world-specific constraints of the game \\
    Language Fluency & Correct grammar, phrasing, and stylistic naturalness \\
    \bottomrule
  \end{tabularx}
  \caption{Qualitative criteria used for evaluating generated quests}
  \label{table:human-eval}
\end{table}

In addition to automated evaluation, a human-based qualitative analysis is conducted on
a representative subset of the generated quests. This assessment focuses on several critical
dimensions: narrative coherence, which examines the logical flow and internal consistency
of the storyline; goal clarity, evaluating whether the objective and associated tasks are
clearly articulated; world consistency, which ensures alignment with the fictional game's
lore and setting; and language fluency, referring to grammatical correctness and natural
phrasing. These qualitative insights help contextualize the numerical results and capture
nuanced attributes of quest quality that are not fully measurable by automated metrics.
The various qualitative criteria used for evaluation are listed in Table~\ref{table:human-eval}.

\newpage
